% GENERATED FILE. Edit canonical lessons, then run npm run generate:books.
% canonical-source-hash: fnv1a64:519242a9f22647b7
% canonical-lessons: ES-C388-empujar, ES-C388-lanzar, ES-C388-borrar, ES-C388-gritar, ES-C388-repaso-cuatro-hedges

\chapter{Push, Throw, Rub Out, Shout}
\label{ch:empujar-y-lanzar}
\hlchaptermodality{388}

\begin{chapteropening}
\textbf{By the end of this chapter:} \emph{I can say empujar, lanzar, borrar and gritar, and name four everyday actions.}

Complete ES-C388-gritar: name four everyday actions, and say one thing you would do with each.
\end{chapteropening}

\section[empujar]{empujar — cousins, not the same word}

\label{lesson:ES-C388-empujar}



\begin{quote}
Say \emph{to fly}, then \emph{to turn}. (\emph{Volar}; \emph{girar}.)
\end{quote}

\subsection*{What to know first}
\begin{itemize}
\raggedright
  \item la puerta — something you push.
  \item la silla — something else you push.
\end{itemize}

\begin{sounds}
\begin{itemize}
\raggedright
  \item \emph{em-pu-JAR}, three beats with the weight on the last. The \emph{j} is the harsh throat sound. $\to$ pronunciation reference
\end{itemize}
\end{sounds}

\begin{cousinweb}
\textbf{empujar} — \textquotedblleft{}to push.\textquotedblright{}

Both dictionaries hedge here, and they hedge in the same direction. The Academy writes \emph{quizá} — perhaps — in front of late Latin \textbf{impulsare}. Corominas writes \emph{probablemente}. Neither will say it outright, and neither has anything better to offer.

They hedge, and there is still hard evidence in the word. Old Spanish spelled it \textbf{empuxar}, with an \emph{x}. That letter is not decoration: an \emph{x} in that spot is what Spanish does to a Latin \textbf{-ls-}. The same change turned Latin \emph{multum} into \emph{mucho}. So whatever the exact Latin verb was, it had an \emph{l} pressed against an \emph{s} inside it — and \emph{impulsare} does.

Behind \emph{impulsare} is \textbf{pellere}, to drive or strike, whose repeating form is \textbf{pulsare}. That is a large family, and English has a great deal of it: \textbf{impulse}, \textbf{impel}, \textbf{propel}, and the \textbf{pulse} in your wrist.

Now the careful part. English \textbf{push} is in that family — but it is \emph{not} the same word as \emph{empujar}. \emph{Push} came through French from the plain \textbf{pulsare}, with nothing on the front. \emph{Empujar} carries an \textbf{in-} that \emph{push} never had. They are cousins who look alike, and calling them one word would be wrong in a way that is easy to check and easy to miss.

Two more traps. The \textbf{pulse} you cook — lentils, chickpeas — is a different word altogether, from Latin \emph{puls}, a thick porridge. And Spanish itself has a \emph{pujar} meaning to bid at an auction, which comes from somewhere else entirely and is no relation to this one.
\end{cousinweb}

\begin{grammarlens}[title={What you've built}]
You can now name the action: \emph{empujar}. And you can read an old spelling as evidence.
\end{grammarlens}

\subsection*{Your turn}
Say these aloud:

\begin{itemize}
\raggedright
  \item \emph{empujar}
  \item \emph{empujar la puerta}
  \item \emph{la silla}
\end{itemize}

\begin{tcolorbox}[breakable,colback=teal!4,colframe=teal!35!black,title={Before you move on}]
What does the \emph{x} in old \emph{empuxar} prove was inside the Latin word? (An \emph{l} pressed against an \emph{s}.) Are \emph{push} and \emph{empujar} the same word? (No — cousins; \emph{empujar} carries an \emph{in-} that \emph{push} never had.)
\end{tcolorbox}
\section[lanzar]{lanzar — a rocket is thrown like a spear}

\label{lesson:ES-C388-lanzar}



\begin{quote}
Say \emph{to turn}, then \emph{to push}. (\emph{Girar}; \emph{empujar}.)
\end{quote}

\subsection*{What to know first}
\begin{itemize}
\raggedright
  \item la mano — what you throw with.
  \item empujar — the gentler neighbour.
\end{itemize}

\begin{sounds}
\begin{itemize}
\raggedright
  \item \emph{lan-ZAR}, two beats with the weight on the second. The \emph{z} is a soft \emph{s} across the Americas and a \emph{th} across much of Spain. $\to$ pronunciation reference
\end{itemize}
\end{sounds}

\begin{cousinweb}
\textbf{lanzar} — \textquotedblleft{}to throw.\textquotedblright{}

Underneath it is a \textbf{lanza}, a lance. Late Latin made the verb \textbf{lanceare}, to handle a lance, and Spanish kept it.

English kept it too, and the English form is the surprise. Around the year 1300 \emph{launchen} meant \textbf{to hurl} — specifically, to hurl the way a spear is hurled. A century later sailors took it for pushing a boat off the beach, and from there it went to rockets, to products, to careers. Every launch in English is still, underneath, a spear leaving a hand.

The plainer relatives are easy: a \textbf{lance}, and a \textbf{lancet}, which is a little lance in a doctor's fingers.

Now where the dictionaries part company, and it is a pattern worth recognising.

Where did Latin get \emph{lancea}? The Academy answers flatly: a word of Celtiberian origin, from the peninsula Spanish is now spoken on. Corominas is far more careful — the Roman writer Varro said it came from Spain, and \emph{quizá}, perhaps, it was Celtiberian. More recent specialists argue for a different ancient Iberian language altogether.

One dictionary printed a possibility as a fact. When a source states plainly what its own sources hedge on, the hedge is the thing to keep.

And one to refuse: the small motor boat called a \textbf{launch} is a different word. It came from Portuguese \emph{lancha}, which the Portuguese took from Malay. It only came to look like the verb afterwards.
\end{cousinweb}

\begin{grammarlens}[title={What you've built}]
You can now name the action: \emph{lanzar}. And you can spot a dictionary sounding surer than its own evidence.
\end{grammarlens}

\subsection*{Your turn}
Say these aloud:

\begin{itemize}
\raggedright
  \item \emph{lanzar}
  \item \emph{empujar y lanzar}
  \item \emph{la mano}
\end{itemize}

\begin{tcolorbox}[breakable,colback=teal!4,colframe=teal!35!black,title={Before you move on}]
What did \emph{launch} mean around 1300? (To hurl, the way a spear is hurled.) Is the boat called a launch the same word? (No — it came from Malay by way of Portuguese.)
\end{tcolorbox}
\section[borrar]{borrar — a verb made out of wool}

\label{lesson:ES-C388-borrar}



\begin{quote}
Say \emph{to push}, then \emph{to throw}. (\emph{Empujar}; \emph{lanzar}.)
\end{quote}

\subsection*{What to know first}
\begin{itemize}
\raggedright
  \item el papel — what you write on.
  \item la clase — where the rubbing out happens.
\end{itemize}

\begin{sounds}
\begin{itemize}
\raggedright
  \item \emph{bo-RRAR}, two beats with the weight on the second. The \emph{rr} is a full roll; the final \emph{r} is a single tap. $\to$ pronunciation reference
\end{itemize}
\end{sounds}

\begin{cousinweb}
\textbf{borrar} — \textquotedblleft{}to erase.\textquotedblright{}

This one is a verb made out of a material. \textbf{Borra} is coarse wool — the rough leftover kind, not the sort you would wear. Late Latin called it \emph{burra}. And a wad of coarse wool was what you rubbed writing off with, back when writing sat on a surface that could be wiped.

So \emph{borrar} does not mean \textquotedblleft{}to erase\textquotedblright{} by any abstract route. It means \emph{to wool-it-off}.

Spanish still keeps the receipt. A rough draft is a \textbf{borrador} — a thing that is \emph{going to be rubbed out}. The same word names the eraser and the board rubber. A first attempt is named for its own impermanence.

English may have a piece of this cloth too, though this is the point at which to slow down.

French had \textbf{bure}, a dark coarse cloth, and its small form \emph{burel} named the cloth spread over a writing table. The cloth gave its name to the table; the table gave its name to the room; the room gave its name to the government department; and the department gave us \textbf{bureaucracy}. That chain is solid all the way along.

The weak joint is at the very start. Dictionaries will not choose between two candidates for where French got \emph{bure} — a Latin word for \emph{red}, or the \emph{burra} above. So the desk really is named after its tablecloth, and whether that tablecloth is \emph{borra}'s relative is not settled. Say the part that holds; leave the part that does not.

One flat refusal: \textbf{borrow} is Old English \emph{borgian}, from the Germanic side of English, and originally meant to stand surety for something. A pure look-alike.
\end{cousinweb}

\begin{grammarlens}[title={What you've built}]
You can now name the action: \emph{borrar}. And you know why a rough draft is called a \emph{borrador}.
\end{grammarlens}

\subsection*{Your turn}
Say these aloud:

\begin{itemize}
\raggedright
  \item \emph{borrar}
  \item \emph{borrar el papel}
  \item \emph{la clase}
\end{itemize}

\begin{tcolorbox}[breakable,colback=teal!4,colframe=teal!35!black,title={Before you move on}]
What was a \emph{borra}? (Coarse wool — what you rubbed writing out with.) Is \emph{borrow} a relative? (No — Old English, with no Latin in it.)
\end{tcolorbox}
\section[gritar]{gritar — a story the Romans told about themselves}

\label{lesson:ES-C388-gritar}



\begin{quote}
Say \emph{to throw}, then \emph{to erase}. (\emph{Lanzar}; \emph{borrar}.)
\end{quote}

\subsection*{What to know first}
\begin{itemize}
\raggedright
  \item la voz — what you shout with.
  \item el llanto — what shouting is not.
\end{itemize}

\begin{sounds}
\begin{itemize}
\raggedright
  \item \emph{gri-TAR}, two beats with the weight on the second. The \emph{g} before \emph{r} is hard, and both \emph{r} sounds are single taps. $\to$ pronunciation reference
\end{itemize}
\end{sounds}

\begin{cousinweb}
\textbf{gritar} — \textquotedblleft{}to shout.\textquotedblright{}

The Academy takes it back to Latin \textbf{quiritare}, to cry out loudly. Corominas agrees that this is probable and calls the origin uncertain even so — which is the more careful of the two positions.

There is a famous story attached, and it is worth telling for the right reason.

Roman writers explained \emph{quiritare} as \emph{calling on the Quirites} — the citizens of Rome — for help. A person in trouble shouts for the townspeople; the verb is named after who you shout at. It is a lovely picture.

It is also almost certainly backwards. Corominas suspects that the ancient dictionary-makers built that gloss out of the resemblance between the two words and then wrote it down as the meaning, which later readers took as evidence. Modern dictionaries set it aside and look elsewhere.

Keep the story, then, but keep it labelled. It is what the Romans believed about their own word, and what a language believes about itself is a real thing to know — as long as you do not file it under fact.

Now the part that is solid. English \textbf{cry} is this same word: Old French \emph{crier}, from the same Latin. And then English did something Spanish did not. \emph{Cry} meant to shout out; over the fifteen-hundreds it drifted across and took over the job of \emph{weep}, until today the shouting sense survives only in phrases like \emph{a far cry} and \emph{hue and cry}.

\emph{Gritar} never made that move. It still means shouting and only shouting — which is why translating English \emph{cry} with it goes wrong so often.

One more to refuse: \textbf{grit} and \textbf{gritty} are Old English \emph{greot}, gravel. Germanic, and no relation.
\end{cousinweb}

\begin{grammarlens}[title={What you've built}]
You can now name the action: \emph{gritar}. And you can hold a story without believing it.
\end{grammarlens}

\subsection*{Your turn}
Say these aloud:

\begin{itemize}
\raggedright
  \item \emph{gritar}
  \item \emph{la voz}
  \item \emph{borrar y gritar}
\end{itemize}

\begin{tcolorbox}[breakable,colback=teal!4,colframe=teal!35!black,title={Before you move on}]
Is the story about calling on the citizens of Rome the real explanation? (No — it looks like a guess the Romans made from the resemblance.) Does \emph{gritar} mean weeping? (No — only shouting; English \emph{cry} moved and Spanish did not.)
\end{tcolorbox}
\section[(review)]{Review — four verbs, four kinds of \textquotedblleft{}related\textquotedblright{}}

\label{lesson:ES-C388-repaso-cuatro-hedges}



\begin{quote}
\textquotedblleft{}I shout.\textquotedblright{} (\emph{Grito}.) \textquotedblleft{}I throw.\textquotedblright{} (\emph{Lanzo}.)
\end{quote}

\begin{grammarlens}[title={What you've built}]
Four ordinary \emph{-ar} verbs. Each one sits next to an English word that looks like it — and each one is related to that English word in a \textbf{different} way.

\noindent
\begin{tabularx}{\linewidth}{@{}>{\raggedright\arraybackslash}X>{\raggedright\arraybackslash}X>{\raggedright\arraybackslash}X@{}}
\toprule
\textbf{Spanish} & \textbf{English lookalike} & \textbf{the truth} \\
\midrule
\emph{lanzar} & launch & \textbf{the same word} \\
\emph{gritar} & cry & \textbf{the same word}, meaning moved \\
\emph{empujar} & push & \textbf{a cousin}, not the same word \\
\emph{borrar} & borrow & \textbf{no relation at all} \\
\bottomrule
\end{tabularx}

Four verbs is a small chapter. Four different answers is a whole method.
\end{grammarlens}

\begin{culture}
\emph{Push} misses being the same word as \emph{empujar} by one piece. Both go back to Latin \emph{pulsare}, but Spanish took the form with \emph{in-} stuck on the front and English took the plain one. Same root, different word.

\emph{Cry} really is \emph{gritar}. Latin \emph{quiritare} went into French and on into English, and then \textbf{English alone} slid the meaning from shouting to weeping. Spanish never did — which is why \emph{gritar} is not the word for tears.

And \emph{borrar} has nothing to do with \emph{borrow}, which is Germanic. \emph{Borrar} comes from \emph{borra}, coarse wool: the pad you rubbed writing out with. A rough draft is still \emph{un borrador} — the thing meant to be rubbed out.

Now the part worth keeping. Look at how the dictionaries actually talk. The Academy writes \emph{quizá} in front of \emph{empujar}; Corominas writes \emph{probablemente}. On \emph{lanzar} they part company outright — the Academy calls its root Celtiberian, Corominas only says \emph{quizá}. On \emph{borrar} neither will choose between two candidates.

A hedge is not a failure. It is a dictionary telling you exactly how much it knows, and a claim with no hedge on it has not been made more certain by leaving the hedge off.
\end{culture}

\subsection*{Your turn}
\begin{itemize}
\raggedright
  \item \emph{Say it:} \emph{empujo} · \emph{lanzo} · \emph{borro} · \emph{grito}
\end{itemize}

\emph{Twice through:}

\begin{tcolorbox}[breakable,colback=teal!4,colframe=teal!35!black,title={Before you move on}]
Does \emph{gritar} mean to weep? (\textbf{No} — English moved that meaning, not Spanish.) What is \emph{un borrador}? (\textbf{A rough draft} — a thing meant to be rubbed out.) When a dictionary writes \emph{quizá}, what is it doing? (\textbf{Telling you how much it knows}.)
\end{tcolorbox}
